% Options for packages loaded elsewhere
\PassOptionsToPackage{unicode}{hyperref}
\PassOptionsToPackage{hyphens}{url}
\PassOptionsToPackage{dvipsnames,svgnames,x11names}{xcolor}
%
\documentclass[
  letterpaper,
  DIV=11,
  numbers=noendperiod]{scrartcl}

\usepackage{amsmath,amssymb}
\usepackage{iftex}
\ifPDFTeX
  \usepackage[T1]{fontenc}
  \usepackage[utf8]{inputenc}
  \usepackage{textcomp} % provide euro and other symbols
\else % if luatex or xetex
  \usepackage{unicode-math}
  \defaultfontfeatures{Scale=MatchLowercase}
  \defaultfontfeatures[\rmfamily]{Ligatures=TeX,Scale=1}
\fi
\usepackage{lmodern}
\ifPDFTeX\else  
    % xetex/luatex font selection
\fi
% Use upquote if available, for straight quotes in verbatim environments
\IfFileExists{upquote.sty}{\usepackage{upquote}}{}
\IfFileExists{microtype.sty}{% use microtype if available
  \usepackage[]{microtype}
  \UseMicrotypeSet[protrusion]{basicmath} % disable protrusion for tt fonts
}{}
\makeatletter
\@ifundefined{KOMAClassName}{% if non-KOMA class
  \IfFileExists{parskip.sty}{%
    \usepackage{parskip}
  }{% else
    \setlength{\parindent}{0pt}
    \setlength{\parskip}{6pt plus 2pt minus 1pt}}
}{% if KOMA class
  \KOMAoptions{parskip=half}}
\makeatother
\usepackage{xcolor}
\setlength{\emergencystretch}{3em} % prevent overfull lines
\setcounter{secnumdepth}{5}
% Make \paragraph and \subparagraph free-standing
\ifx\paragraph\undefined\else
  \let\oldparagraph\paragraph
  \renewcommand{\paragraph}[1]{\oldparagraph{#1}\mbox{}}
\fi
\ifx\subparagraph\undefined\else
  \let\oldsubparagraph\subparagraph
  \renewcommand{\subparagraph}[1]{\oldsubparagraph{#1}\mbox{}}
\fi

\usepackage{color}
\usepackage{fancyvrb}
\newcommand{\VerbBar}{|}
\newcommand{\VERB}{\Verb[commandchars=\\\{\}]}
\DefineVerbatimEnvironment{Highlighting}{Verbatim}{commandchars=\\\{\}}
% Add ',fontsize=\small' for more characters per line
\usepackage{framed}
\definecolor{shadecolor}{RGB}{241,243,245}
\newenvironment{Shaded}{\begin{snugshade}}{\end{snugshade}}
\newcommand{\AlertTok}[1]{\textcolor[rgb]{0.68,0.00,0.00}{#1}}
\newcommand{\AnnotationTok}[1]{\textcolor[rgb]{0.37,0.37,0.37}{#1}}
\newcommand{\AttributeTok}[1]{\textcolor[rgb]{0.40,0.45,0.13}{#1}}
\newcommand{\BaseNTok}[1]{\textcolor[rgb]{0.68,0.00,0.00}{#1}}
\newcommand{\BuiltInTok}[1]{\textcolor[rgb]{0.00,0.23,0.31}{#1}}
\newcommand{\CharTok}[1]{\textcolor[rgb]{0.13,0.47,0.30}{#1}}
\newcommand{\CommentTok}[1]{\textcolor[rgb]{0.37,0.37,0.37}{#1}}
\newcommand{\CommentVarTok}[1]{\textcolor[rgb]{0.37,0.37,0.37}{\textit{#1}}}
\newcommand{\ConstantTok}[1]{\textcolor[rgb]{0.56,0.35,0.01}{#1}}
\newcommand{\ControlFlowTok}[1]{\textcolor[rgb]{0.00,0.23,0.31}{#1}}
\newcommand{\DataTypeTok}[1]{\textcolor[rgb]{0.68,0.00,0.00}{#1}}
\newcommand{\DecValTok}[1]{\textcolor[rgb]{0.68,0.00,0.00}{#1}}
\newcommand{\DocumentationTok}[1]{\textcolor[rgb]{0.37,0.37,0.37}{\textit{#1}}}
\newcommand{\ErrorTok}[1]{\textcolor[rgb]{0.68,0.00,0.00}{#1}}
\newcommand{\ExtensionTok}[1]{\textcolor[rgb]{0.00,0.23,0.31}{#1}}
\newcommand{\FloatTok}[1]{\textcolor[rgb]{0.68,0.00,0.00}{#1}}
\newcommand{\FunctionTok}[1]{\textcolor[rgb]{0.28,0.35,0.67}{#1}}
\newcommand{\ImportTok}[1]{\textcolor[rgb]{0.00,0.46,0.62}{#1}}
\newcommand{\InformationTok}[1]{\textcolor[rgb]{0.37,0.37,0.37}{#1}}
\newcommand{\KeywordTok}[1]{\textcolor[rgb]{0.00,0.23,0.31}{#1}}
\newcommand{\NormalTok}[1]{\textcolor[rgb]{0.00,0.23,0.31}{#1}}
\newcommand{\OperatorTok}[1]{\textcolor[rgb]{0.37,0.37,0.37}{#1}}
\newcommand{\OtherTok}[1]{\textcolor[rgb]{0.00,0.23,0.31}{#1}}
\newcommand{\PreprocessorTok}[1]{\textcolor[rgb]{0.68,0.00,0.00}{#1}}
\newcommand{\RegionMarkerTok}[1]{\textcolor[rgb]{0.00,0.23,0.31}{#1}}
\newcommand{\SpecialCharTok}[1]{\textcolor[rgb]{0.37,0.37,0.37}{#1}}
\newcommand{\SpecialStringTok}[1]{\textcolor[rgb]{0.13,0.47,0.30}{#1}}
\newcommand{\StringTok}[1]{\textcolor[rgb]{0.13,0.47,0.30}{#1}}
\newcommand{\VariableTok}[1]{\textcolor[rgb]{0.07,0.07,0.07}{#1}}
\newcommand{\VerbatimStringTok}[1]{\textcolor[rgb]{0.13,0.47,0.30}{#1}}
\newcommand{\WarningTok}[1]{\textcolor[rgb]{0.37,0.37,0.37}{\textit{#1}}}

\providecommand{\tightlist}{%
  \setlength{\itemsep}{0pt}\setlength{\parskip}{0pt}}\usepackage{longtable,booktabs,array}
\usepackage{calc} % for calculating minipage widths
% Correct order of tables after \paragraph or \subparagraph
\usepackage{etoolbox}
\makeatletter
\patchcmd\longtable{\par}{\if@noskipsec\mbox{}\fi\par}{}{}
\makeatother
% Allow footnotes in longtable head/foot
\IfFileExists{footnotehyper.sty}{\usepackage{footnotehyper}}{\usepackage{footnote}}
\makesavenoteenv{longtable}
\usepackage{graphicx}
\makeatletter
\def\maxwidth{\ifdim\Gin@nat@width>\linewidth\linewidth\else\Gin@nat@width\fi}
\def\maxheight{\ifdim\Gin@nat@height>\textheight\textheight\else\Gin@nat@height\fi}
\makeatother
% Scale images if necessary, so that they will not overflow the page
% margins by default, and it is still possible to overwrite the defaults
% using explicit options in \includegraphics[width, height, ...]{}
\setkeys{Gin}{width=\maxwidth,height=\maxheight,keepaspectratio}
% Set default figure placement to htbp
\makeatletter
\def\fps@figure{htbp}
\makeatother

\KOMAoption{captions}{tableheading}
\makeatletter
\makeatother
\makeatletter
\makeatother
\makeatletter
\@ifpackageloaded{caption}{}{\usepackage{caption}}
\AtBeginDocument{%
\ifdefined\contentsname
  \renewcommand*\contentsname{Table of contents}
\else
  \newcommand\contentsname{Table of contents}
\fi
\ifdefined\listfigurename
  \renewcommand*\listfigurename{List of Figures}
\else
  \newcommand\listfigurename{List of Figures}
\fi
\ifdefined\listtablename
  \renewcommand*\listtablename{List of Tables}
\else
  \newcommand\listtablename{List of Tables}
\fi
\ifdefined\figurename
  \renewcommand*\figurename{Figure}
\else
  \newcommand\figurename{Figure}
\fi
\ifdefined\tablename
  \renewcommand*\tablename{Table}
\else
  \newcommand\tablename{Table}
\fi
}
\@ifpackageloaded{float}{}{\usepackage{float}}
\floatstyle{ruled}
\@ifundefined{c@chapter}{\newfloat{codelisting}{h}{lop}}{\newfloat{codelisting}{h}{lop}[chapter]}
\floatname{codelisting}{Listing}
\newcommand*\listoflistings{\listof{codelisting}{List of Listings}}
\makeatother
\makeatletter
\@ifpackageloaded{caption}{}{\usepackage{caption}}
\@ifpackageloaded{subcaption}{}{\usepackage{subcaption}}
\makeatother
\makeatletter
\@ifpackageloaded{tcolorbox}{}{\usepackage[skins,breakable]{tcolorbox}}
\makeatother
\makeatletter
\@ifundefined{shadecolor}{\definecolor{shadecolor}{rgb}{.97, .97, .97}}
\makeatother
\makeatletter
\makeatother
\makeatletter
\makeatother
\ifLuaTeX
  \usepackage{selnolig}  % disable illegal ligatures
\fi
\IfFileExists{bookmark.sty}{\usepackage{bookmark}}{\usepackage{hyperref}}
\IfFileExists{xurl.sty}{\usepackage{xurl}}{} % add URL line breaks if available
\urlstyle{same} % disable monospaced font for URLs
\hypersetup{
  pdftitle={Reproducing WLasso on breast cancer gene expression data},
  pdfauthor={Mariana Garcia Mejia},
  colorlinks=true,
  linkcolor={blue},
  filecolor={Maroon},
  citecolor={Blue},
  urlcolor={Blue},
  pdfcreator={LaTeX via pandoc}}

\title{Reproducing WLasso on breast cancer gene expression data}
\author{Mariana Garcia Mejia}
\date{October 22, 2026}

\begin{document}
\maketitle
\ifdefined\Shaded\renewenvironment{Shaded}{\begin{tcolorbox}[interior hidden, borderline west={3pt}{0pt}{shadecolor}, enhanced, sharp corners, boxrule=0pt, frame hidden, breakable]}{\end{tcolorbox}}\fi

\hypertarget{reproducing-wlasso-on-breast-cancer-gene-expression-data-gse2990}{%
\subsection{Reproducing WLasso on breast cancer gene expression data
(GSE2990)}\label{reproducing-wlasso-on-breast-cancer-gene-expression-data-gse2990}}

\hypertarget{data-and-preprocessing}{%
\subsubsection{Data and preprocessing}\label{data-and-preprocessing}}

We use data from the Gene Expression Omnibus
\href{https://www.ncbi.nlm.nih.gov/geo/}{Gene Expression Omnibus} with
GEO Accession GSE2990 as used by the authors (Zhu et al., 2020). The
data contains \(n = 189\) tumours and 22,283 probes (HG-U133A array).
The values range from 2.66 to 15.13, which is consistent with
RMA-normalised data.

\begin{Shaded}
\begin{Highlighting}[]
\CommentTok{\# Load the gene data from Gene Expression Omnibus}
\NormalTok{expr }\OtherTok{\textless{}{-}} \FunctionTok{read.delim}\NormalTok{(}\StringTok{"GSE2990\_series\_matrix.txt"}\NormalTok{,}
                   \AttributeTok{comment.char =} \StringTok{"!"}\NormalTok{, }\AttributeTok{row.names =} \DecValTok{1}\NormalTok{)}
\FunctionTok{dim}\NormalTok{(expr)}
\end{Highlighting}
\end{Shaded}

\begin{verbatim}
[1] 22283   189
\end{verbatim}

\begin{Shaded}
\begin{Highlighting}[]
\CommentTok{\# Values appear to be normalized (between 2 and 15).}
\FunctionTok{range}\NormalTok{(}\FunctionTok{as.matrix}\NormalTok{(expr), }\AttributeTok{na.rm =} \ConstantTok{TRUE}\NormalTok{)}
\end{Highlighting}
\end{Shaded}

\begin{verbatim}
[1]  2.656036 15.129603
\end{verbatim}

As in Zhu et al.~(2021), we kept probes with IQR \(\geq 1.5\) (1,155
probes) and removed probes without a gene annotation (49 removed, 1,106
left).

\begin{Shaded}
\begin{Highlighting}[]
\CommentTok{\# First cleaning step: remove probes with IQR \textless{} 1.5}
\NormalTok{iqr\_values }\OtherTok{\textless{}{-}} \FunctionTok{apply}\NormalTok{(expr, }\DecValTok{1}\NormalTok{, IQR)}
\CommentTok{\# summary(iqr\_values)}
\NormalTok{keep }\OtherTok{\textless{}{-}}\NormalTok{ iqr\_values }\SpecialCharTok{\textgreater{}=} \FloatTok{1.5}
\CommentTok{\# sum(keep)}
\NormalTok{expr\_filt }\OtherTok{\textless{}{-}}\NormalTok{ expr[keep, ]}
\CommentTok{\# dim(expr\_filt)}

\CommentTok{\# Second step: Remove genes without gene annotation}
\CommentTok{\# BiocManager::install("hgu133a.db")}
\FunctionTok{library}\NormalTok{(hgu133a.db)}

\NormalTok{probe\_ids }\OtherTok{\textless{}{-}} \FunctionTok{rownames}\NormalTok{(expr\_filt)}
\NormalTok{symbols }\OtherTok{\textless{}{-}} \FunctionTok{mapIds}\NormalTok{(hgu133a.db,}
                  \AttributeTok{keys =}\NormalTok{ probe\_ids,}
                  \AttributeTok{column =} \StringTok{"SYMBOL"}\NormalTok{,}
                  \AttributeTok{keytype =} \StringTok{"PROBEID"}\NormalTok{,}
                  \AttributeTok{multiVals =} \StringTok{"first"}\NormalTok{)}

\CommentTok{\# sum(is.na(symbols))}

\CommentTok{\# Remove}
\NormalTok{has\_symbol }\OtherTok{\textless{}{-}} \SpecialCharTok{!}\FunctionTok{is.na}\NormalTok{(symbols)}
\NormalTok{expr\_ann }\OtherTok{\textless{}{-}}\NormalTok{ expr\_filt[has\_symbol, ]}
\NormalTok{symbols\_ann }\OtherTok{\textless{}{-}}\NormalTok{ symbols[has\_symbol]}

\FunctionTok{dim}\NormalTok{(expr\_ann)}
\end{Highlighting}
\end{Shaded}

\begin{verbatim}
[1] 1106  189
\end{verbatim}

The ESR1 probe is a gene associated with breast cancer. We use it as the
response \(y\) and the remaining 1,105 probes as predictors \(X\). As
with regular As is convention in Lasso, we standardised each predictor.
We also centred \(y\), because the criterion has no intercept. The paper
reports \(p = 1{,}111\), so our probe set differs slightly,
\(p = 1{,}106\).

\begin{Shaded}
\begin{Highlighting}[]
\CommentTok{\# Use ESR1 gene as the gene to predict.}
\NormalTok{all\_symbols }\OtherTok{\textless{}{-}} \FunctionTok{mapIds}\NormalTok{(hgu133a.db,}
                      \AttributeTok{keys =} \FunctionTok{rownames}\NormalTok{(expr),}
                      \AttributeTok{column =} \StringTok{"SYMBOL"}\NormalTok{,}
                      \AttributeTok{keytype =} \StringTok{"PROBEID"}\NormalTok{,}
                      \AttributeTok{multiVals =} \StringTok{"first"}\NormalTok{)}

\NormalTok{esr1\_probes }\OtherTok{\textless{}{-}} \FunctionTok{names}\NormalTok{(all\_symbols)[all\_symbols }\SpecialCharTok{\%in\%} \StringTok{"ESR1"}\NormalTok{]}
\CommentTok{\# esr1\_probes}
\CommentTok{\# esr1\_probes \%in\% rownames(expr\_ann)}

\CommentTok{\# Pull the ESR Variable left after IQR + Annotation filter and use as y}
\NormalTok{esr1\_probe }\OtherTok{\textless{}{-}}\NormalTok{ esr1\_probes[}\DecValTok{1}\NormalTok{]}
\CommentTok{\# esr1\_probe}
\NormalTok{y }\OtherTok{\textless{}{-}} \FunctionTok{as.numeric}\NormalTok{(expr\_ann[esr1\_probe, ])}
\NormalTok{expr\_X }\OtherTok{\textless{}{-}}\NormalTok{ expr\_ann[}\FunctionTok{rownames}\NormalTok{(expr\_ann) }\SpecialCharTok{!=}\NormalTok{ esr1\_probe, ]}

\CommentTok{\# length(y)}
\CommentTok{\# dim(expr\_X)}

\CommentTok{\# Standardize X.}
\NormalTok{X }\OtherTok{\textless{}{-}} \FunctionTok{scale}\NormalTok{(}\FunctionTok{t}\NormalTok{(expr\_X))}
\CommentTok{\# dim(X)}
\CommentTok{\# round(colMeans(X)[1:3], 6)}
\CommentTok{\# apply(X, 2, sd)[1:3]}

\CommentTok{\# Mean{-}center Y.}
\NormalTok{y\_c }\OtherTok{\textless{}{-}}\NormalTok{ y }\SpecialCharTok{{-}} \FunctionTok{mean}\NormalTok{(y)}
\CommentTok{\# mean(y\_c) \# should be close to zero.}
\end{Highlighting}
\end{Shaded}

\hypertarget{what-the-gene-annotation-does}{%
\subsubsection{What the gene annotation
does}\label{what-the-gene-annotation-does}}

Each column \(j \in \{1,\dots,p\}\) of \(X\) is a probe, and the array
itself only knows probe IDs. An annotation is a lookup map \(g\) from
probe indices to gene symbols (or to \(\mathrm{NA}\) when the probe has
no known gene). We used the Bioconductor package \texttt{hgu133a.db} for
this array. It never enters the estimation of \(\beta\), only the steps
before and after:

\begin{itemize}
\tightlist
\item
  \textbf{Before:} we drop probes with \(g(j) = \mathrm{NA}\), and we
  find the probe \(j^*\) with \(g(j^*) = \text{ESR1}\) to build
  \(y = X_{j^*}\) and take it out of the predictors.
\item
  \textbf{After:} the method returns a support
  \(\hat S = \{ j : \hat\beta_j \neq 0 \}\) of probes. The map turns it
  into a list of genes \(g(\hat S)\). Several probes can map to the same
  gene, so \(|g(\hat S)| \leq |\hat S|\).
\end{itemize}

\hypertarget{method}{%
\subsubsection{Method}\label{method}}

The model is \(y = X\beta + \epsilon\) with
\(\mathrm{Cor}(X) = \Sigma\). WLasso whitens the predictors with
\(\tilde{X} = X\Sigma^{-1/2}\), so that
\(y = \tilde{X}\tilde\beta + \epsilon\) with
\(\tilde\beta = \Sigma^{1/2}\beta\), and the error term is unchanged. It
then minimises the generalised Lasso criterion
\[\|y - \tilde{X}\tilde\beta\|_2^2 + \lambda \|\Sigma^{-1/2}\tilde\beta\|_1,\]
and thresholds the estimates to the top \(K\) (for \(\tilde\beta\)) and
top \(M\) (for \(\beta\)) coefficients, with \(K\) and \(M\) chosen by
an MSE-ratio rule controlled by \(\gamma\). \(\Sigma\) was estimated
with the two-block structure of Section 2.5:
\((\hat\alpha_1,\hat\alpha_2,\hat\alpha_3) = (0.19, 0.27, 0.55)\),
against \((0.17, 0.21, 0.52)\) in the paper.

\hypertarget{a-problem-with-the-package-and-our-patch}{%
\subsubsection{A problem with the package and our
patch}\label{a-problem-with-the-package-and-our-patch}}

With the \texttt{WLasso} package (v1.0), \texttt{Whitening\_Lasso()}
returned an empty model on these data (the lowest MSE was at the first,
all-zero step), although a plain Lasso on the same data selected 62
probes. The package works on its own example data and on simulated data
of our size (\(n = 189\), \(p = 1105\)). Reading the source, in the step
that chooses \(M\), the MSE is computed from the coefficients sorted by
absolute value, not from the coefficients in their original order, so
each coefficient multiplies the wrong column of \(X\). We think this is
a bug that does little harm when all correlations are positive, as in
the paper's simulations. We changed that one argument in a copy of the
function (not an official fix, and not confirmed with the authors), and
results below come from that modified copy.

\begin{Shaded}
\begin{Highlighting}[]
\CommentTok{\# Install WLasso package}
\CommentTok{\# install.packages("WLasso")}
\FunctionTok{library}\NormalTok{(WLasso)}
\CommentTok{\# help(package = "WLasso")}
\end{Highlighting}
\end{Shaded}

\begin{Shaded}
\begin{Highlighting}[]
\CommentTok{\# Step 1: Sigma Estimation}
\NormalTok{Sigma\_est }\OtherTok{\textless{}{-}} \FunctionTok{Sigma\_Estimation}\NormalTok{(X)}
\FunctionTok{round}\NormalTok{(Sigma\_est}\SpecialCharTok{$}\NormalTok{alpha, }\DecValTok{2}\NormalTok{) }\CommentTok{\# paper reports 0.17, 0.21 and 0.52, we get slightly different}
\end{Highlighting}
\end{Shaded}

\begin{verbatim}
[1] 0.19 0.27 0.55
\end{verbatim}

\begin{Shaded}
\begin{Highlighting}[]
\CommentTok{\# We get: 0.19 0.27 0.55}

\CommentTok{\# Number of probes the clustering put into one of the two blocks}
\CommentTok{\# We get 181 out of 1105}
\FunctionTok{length}\NormalTok{(Sigma\_est}\SpecialCharTok{$}\NormalTok{group\_act)}
\end{Highlighting}
\end{Shaded}

\begin{verbatim}
[1] 181
\end{verbatim}

\begin{Shaded}
\begin{Highlighting}[]
\CommentTok{\# Step 2: Model transformation}
\NormalTok{Sigma\_estmat }\OtherTok{\textless{}{-}}\NormalTok{ Sigma\_est}\SpecialCharTok{$}\NormalTok{mat}
\end{Highlighting}
\end{Shaded}

\begin{Shaded}
\begin{Highlighting}[]
\CommentTok{\# Fix bug on the paper}
\NormalTok{src }\OtherTok{\textless{}{-}} \FunctionTok{paste}\NormalTok{(}\FunctionTok{deparse}\NormalTok{(WLasso}\SpecialCharTok{::}\NormalTok{Whitening\_Lasso), }\AttributeTok{collapse =} \StringTok{"}\SpecialCharTok{\textbackslash{}n}\StringTok{"}\NormalTok{)}
\NormalTok{src2 }\OtherTok{\textless{}{-}} \FunctionTok{gsub}\NormalTok{(}\StringTok{"x}\SpecialCharTok{\textbackslash{}\textbackslash{}}\StringTok{s*=}\SpecialCharTok{\textbackslash{}\textbackslash{}}\StringTok{s*beta\_sort,}\SpecialCharTok{\textbackslash{}\textbackslash{}}\StringTok{s*sorted\_vect}\SpecialCharTok{\textbackslash{}\textbackslash{}}\StringTok{s*=}\SpecialCharTok{\textbackslash{}\textbackslash{}}\StringTok{s*beta\_sort"}\NormalTok{,}
             \StringTok{"x = beta\_final0, sorted\_vect = beta\_sort"}\NormalTok{, src)}

\CommentTok{\# identical(src, src2)   \# we want FALSE this time}

\NormalTok{WL\_fixed }\OtherTok{\textless{}{-}} \FunctionTok{eval}\NormalTok{(}\FunctionTok{parse}\NormalTok{(}\AttributeTok{text =}\NormalTok{ src2))}
\FunctionTok{environment}\NormalTok{(WL\_fixed) }\OtherTok{\textless{}{-}} \FunctionTok{asNamespace}\NormalTok{(}\StringTok{"WLasso"}\NormalTok{)}
\CommentTok{\# grep("beta\_sort", strsplit(src, "\textbackslash{}n")[[1]], value = TRUE)}
\CommentTok{\# grep("beta\_interm", strsplit(src2, "\textbackslash{}n")[[1]], value = TRUE)}
\end{Highlighting}
\end{Shaded}

\begin{Shaded}
\begin{Highlighting}[]
\CommentTok{\# Fit the model}
\CommentTok{\# mod\_fix \textless{}{-} WL\_fixed(X = X, Y = y\_c, Sigma = Sigma\_estmat,}
                    \CommentTok{\# gamma = 0.95, maxsteps = 500)}

\CommentTok{\# range(mod\_fix$mse)}
\CommentTok{\# sum(mod\_fix$beta.min != 0)}
\end{Highlighting}
\end{Shaded}

\begin{Shaded}
\begin{Highlighting}[]
\CommentTok{\# We also fit regular Lasso for comparison:}
\CommentTok{\#install.packages("glmnet")}
\FunctionTok{library}\NormalTok{(glmnet)}
\end{Highlighting}
\end{Shaded}

\begin{verbatim}
Warning: package 'glmnet' was built under R version 4.3.3
\end{verbatim}

\begin{verbatim}
Loading required package: Matrix
\end{verbatim}

\begin{verbatim}

Attaching package: 'Matrix'
\end{verbatim}

\begin{verbatim}
The following object is masked from 'package:S4Vectors':

    expand
\end{verbatim}

\begin{verbatim}
Loaded glmnet 4.1-9
\end{verbatim}

\begin{Shaded}
\begin{Highlighting}[]
\NormalTok{cvfit }\OtherTok{\textless{}{-}} \FunctionTok{cv.glmnet}\NormalTok{(X, y\_c)}
\FunctionTok{sum}\NormalTok{(}\FunctionTok{coef}\NormalTok{(cvfit, }\AttributeTok{s =} \StringTok{"lambda.min"}\NormalTok{)[}\SpecialCharTok{{-}}\DecValTok{1}\NormalTok{] }\SpecialCharTok{!=} \DecValTok{0}\NormalTok{)}
\end{Highlighting}
\end{Shaded}

\begin{verbatim}
[1] 60
\end{verbatim}

\begin{Shaded}
\begin{Highlighting}[]
\CommentTok{\# returns 61 coefficients!!!!}
\end{Highlighting}
\end{Shaded}

\hypertarget{results}{%
\subsubsection{Results}\label{results}}

\begin{longtable}[]{@{}lll@{}}
\toprule\noalign{}
Method & Probes & Genes \\
\midrule\noalign{}
\endhead
\bottomrule\noalign{}
\endlastfoot
WLasso (\(\gamma = 0.95\), patched) & 59 & 56 \\
Lasso (cross-validation) & 62 & 59 \\
\end{longtable}

The paper reports 63 and 66 genes. The two methods share 38 genes in our
run, against 8 in the paper. NAT1 and GATA3, both discussed in the
paper, are selected by both methods. Varying \(\gamma\) changes the size
of the selection (42, 56 and 72 genes for
\(\gamma = 0.90, 0.95, 0.99\)), but 36 genes are selected at all three
values. Some genes discussed in the paper are unstable: TOP2A is
selected at 0.90 and 0.99 but not at 0.95, and CXCL12 only at 0.99.

\hypertarget{limitations}{%
\subsubsection{Limitations}\label{limitations}}

The true set of relevant genes is unknown, so this is an illustration,
as in the paper, and cannot show that one method is better. We only
checked stability with respect to \(\gamma\), not \(\hat\Sigma\) or the
sample. The results rely on the patched function and the current
annotation version.

\textbf{Reference.} Zhu, W., Lévy-Leduc, C. and Ternès, N. (2021). A
variable selection approach for highly correlated predictors in
high-dimensional genomic data. \emph{Bioinformatics}, 37(16),
2238--2244.



\end{document}
